\SourceLeaf{470}
\par\Needspace{9\baselineskip}\bigskip
\pdfbookmark[3]{Exposition, 20}{pallas-II-article20}
\begin{center}20.\end{center}
\begingroup
\setlength{\lineskip}{4pt}\setlength{\lineskiplimit}{4pt}
\noindent\hspace*{\parindent}Les valeurs variables des élémens, qu’on obtient par cette méthode, seront justes aux quantités près de l’ordre du quarré des masses perturbatrices. Si on désire une précision encore plus grande, on n’a qu’à refaire ce même calcul sur les élémens variables qu’on vient de trouver: de cette manière on aura des résultats justes aux quantités près du troisième ordre. Il serait superflu de répéter le calcul encore une fois: on pourra même se permettre dans le second calcul, de supposer. les élémens constans pendant quelque tems, c. à d. d’employer dans le calcul des variations instantanées un même système d’élémens pour plusieurs époques antérieures ou postérieures à celle à laquelle il répond. Généralement on ne changera pas de système, avant que la différence entre ce système et celui qu’il faudrait lui substituer soit devenue assés sensible: par conséquent, plus les variations instantanées des élémens sont petites, plus le nombre des époques sera grand, pour lesquelles il est permis de se servir du même système d’élémens.
\par Il sera un peu plus commode pour la pratique, de multiplier les valeurs de \(\dfrac{\mathrm{d}i}{\mathrm{d}t},\ \dfrac{\mathrm{d}\AscendingNode}{\mathrm{d}t},\ \dfrac{\mathrm{d}\varphi}{\mathrm{d}t},\ \dfrac{\mathrm{d}\tilde\omega}{\mathrm{d}t},\ \dfrac{\mathrm{d}\epsilon}{\mathrm{d}t}\) par \(\theta\) et celles de \(\dfrac{\mathrm{d}n}{\mathrm{d}t}\) par \(\theta\theta\), avant de former les séries des différences et des sommes: il est superflu, de remarquer qu’il convient d’exprimer d’abord ces quantités en secondes. La grandeur la plus convenables de l’intervalle \(\theta\) dépend en partie de la vitesse du mouvement des planètes perturbante et perturbée, et en partie du degré de précision qu’on veut atteindre. L’auteur calculant d’après cette méthode les perturbations de Pallas produites par Jupiter pendant les années \mbox{1803 . . . . 1816} a choisi l’intervalle de 50 jours: c’est au calculateur intelligent de juger, si, dans d’autres circonstances, il conviendra de diminuer ou d’augmenter cet intervalle. Nous devons supprimer ici plusieurs autres remarques pratiques, qui s’offriront d’elles-mêmes à tout calculateur exercé, qui exécutera des travaux semblables.
\par\endgroup
\begin{center}\rule{0.18\textwidth}{0.35pt}\end{center}
